\documentclass[../../../include/open-logic-section]{subfiles}

\begin{document}

\olfileid[pt]{sth}{cardinals}{classing}
\olsection{Finito, Enumerável, Não Enumerável}

Agora que fomos introduzidos aos cardinais, vale a pena dedicar um
pouco de tempo a variedades de cardinais;
nomeadamente, cardinais finitos, !!{enumerable} e !!{nonenumerable}.

Os nossos dois primeiros resultados implicam que os cardinais finitos serão exatamente
os ordinais finitos, que definimos como os nossos \emph{números naturais}
em \olref[z][infinity-again]{defnomega}:

\begin{prop}\ollabel{finitecardisoequal}
Sejam $n, m \in \omega$. Então $n = m$ sse $\cardeq{n}{m}$.
\end{prop}

\begin{proof}
O sentido \emph{da esquerda para a direita} é trivial. Para provar \emph{da direita para a esquerda},
suponha-se $\cardeq{n}{m}$ embora $n \neq m$. Por tricotomia, ou $n
\in m$ ou $m \in n$; suponha-se, sem perda de generalidade, $n \in m$.
Então $n \subsetneq m$ e existe !!a{bijection} $f \colon m \to n$,
logo $m$ é infinito de Dedekind, o que contradiz
\olref[z][infinity-again]{naturalnumbersarentinfinite}.
\end{proof}

\begin{cor}\ollabel{naturalsarecardinals}
Se $n \in \omega$, então $n$ é um cardinal.
\end{cor}

\begin{proof}
Imediato.
\end{proof}
\noindent
Segue-se também que várias noções razoáveis do que possa significar
descrever um cardinal como ``finito'' ou ``infinito'' coincidem:
\begin{thm}\ollabel{generalinfinitycharacter}Para qualquer conjunto $A$, as seguintes afirmações são equivalentes:
\begin{enumerate}
	\item\ollabel{card:notinomega} $\card{A} \notin \omega$, isto é,
	$A$ não é um número natural;
	\item\ollabel{card:omegaplus} $\omega \leq \card{A}$;
	\item\ollabel{card:infinite} $A$ é infinito de Dedekind.
\end{enumerate}
\end{thm}

\begin{proof}
Por \olref[ord-arithmetic][using-addition]{ordinfinitycharacter},
\olref[cardinals][cardsasords]{lem:CardinalsBehaveRight}, e
\olref{naturalsarecardinals}.
\end{proof}

Isto autoriza a seguinte \emph{definição} de algumas noções que usamos
de forma informal na \olref[sfr][][]{part}:

\begin{defn}\ollabel{defnfinite}
Diz-se que $A$ é \emph{finito} sse $\card{A}$ é um número natural,
isto é, $\card{A} \in \omega$. Caso contrário, diz-se que $A$ é
\emph{infinito}.
\end{defn}
\noindent
Note-se que esta definição é apresentada no contexto de
$\ZFC$. Afinal, precisamos do Bom ordenamento para garantir que todo o conjunto
tem uma cardinalidade. E, de fato, sem o Bom ordenamento, pode existir um
conjunto que nem é finito nem infinito de Dedekind. Voltaremos a
este tipo de questão no \olref[choice][]{chap}. Por ora, continuamos a
assumir o Bom ordenamento.

Passemos agora dos cardinais finitos aos cardinais infinitos.
Eis dois pontos elementares:

\begin{cor}\ollabel{omegaisacardinal}
$\omega$ é o menor cardinal infinito.
\end{cor}

\begin{proof}
$\omega$ é um cardinal, pois $\omega$ é infinito de Dedekind e se
$\cardeq{\omega}{n}$ para algum $n \in \omega$ então $n$ seria infinito de Dedekind,
contrariando
\olref[z][infinity-again]{naturalnumbersarentinfinite}. Ora,
$\omega$ é o menor cardinal infinito por definição.
\end{proof}

\begin{cor}
Todo o cardinal infinito é um ordinal limite.
\end{cor}

\begin{proof}
Seja $\alpha$ um ordinal sucessor infinito, logo $\alpha = \beta
\ordplus 1$ para algum $\beta$. Por \olref{finitecardisoequal}, $\beta$
também é infinito, logo $\cardeq{\beta}{\beta \ordplus 1}$ por
\olref[ord-arithmetic][using-addition]{ordinfinitycharacter}. Ora,
$\card{\beta} = \card{\beta\ordplus 1} = \card{\alpha}$ por
\olref[cardinals][cardsasords]{lem:CardinalsBehaveRight}, logo
$\alpha \neq \card{\alpha}$.
\end{proof}

Já na \olref[sfr][siz][enm-alt]{defn:enumerable}, assinalamos que
podemos distinguir conjuntos infinitos !!{enumerable} de conjuntos infinitos
!!{nonenumerable}. Essa definição conduz naturalmente ao seguinte:

\begin{prop}
$A$ é !!{enumerable} sse $\card{A} \leq \omega$, e $A$ é
!!{nonenumerable} sse $\omega < \card{A}$.
\end{prop}

\begin{proof}
Por tricotomia, as duas afirmações são equivalentes, bastando provar
que $A$ é !!{enumerable} sse $\card{A} \leq \omega$. Para
\emph{da direita para a esquerda}: se $\card{A} \leq \omega$, então
$\cardle{A}{\omega}$ por
\olref[cardinals][cardsasords]{lem:CardinalsBehaveRight} e
\olref{omegaisacardinal}. Para \emph{da esquerda para a direita}: suponha-se $A$
!!{enumerable}; então, por \olref[sfr][siz][enm-alt]{defn:enumerable}, há
três casos possíveis:
\begin{enumerate}
	\item se $A = \emptyset$, então $\card{A} = 0 \in \omega$, por
	\olref{naturalsarecardinals} e
	\olref[cardinals][cardsasords]{lem:CardinalsBehaveRight}.
	\item se $\cardeq{n}{A}$, então $\card{A} = n \in \omega$, por
	\olref{naturalsarecardinals} e
	\olref[cardinals][cardsasords]{lem:CardinalsBehaveRight}.
	\item se $\cardeq{\omega}{A}$, então $\card{A} = \omega$, por \olref{omegaisacardinal}.
\end{enumerate}
Em todos os casos, $\card{A} \leq \omega$.
\end{proof}
\noindent
De fato, $\omega$ tem um lugar especial. Embora existam muitos ordinais contáveis:

\begin{cor}
$\omega$ é o único cardinal infinito !!{enumerable}.
\end{cor}

\begin{proof}
Seja $\cardfont{a}$ um cardinal infinito enumerável. Como
$\cardfont{a}$ é infinito, $\omega \leq \cardfont{a}$. Como
$\cardfont{a}$ é um cardinal infinito enumerável, $\cardfont{a} =
\card{\cardfont{a}} \leq \omega$. Logo $\cardfont{a} = \omega$ por
tricotomia. \end{proof}

Claro que existem infinitos cardinais. Poderíamos perguntar:
\emph{Quantos cardinais há?} Os resultados seguintes mostram que
convém reconsiderar essa pergunta.

\begin{prop}\ollabel{unioncardinalscardinal}
Se todo o membro de $X$ é um cardinal, então $\bigcup X$ é um cardinal.
\end{prop}

\begin{proof}
É fácil verificar que $\bigcup X$ é um ordinal. Seja $\alpha \in
\bigcup X$ um ordinal; então $\alpha \in \cardfont{b} \in X$ para
algum cardinal $\cardfont{b}$. Como $\cardfont{b}$ é um cardinal,
$\cardless{\alpha}{\cardfont{b}}$. Como $\cardfont{b} \subseteq
\bigcup X$, tem-se $\cardle{\cardfont{b}}{\bigcup X}$, e portanto
$\cardneq{\alpha}{\bigcup X}$. Generalizando, $\bigcup X$ é um
cardinal.
\end{proof}

\begin{thm}\ollabel{lem:NoLargestCardinal}
Não existe um cardinal máximo.
\end{thm}

\begin{proof}
Para qualquer cardinal $\cardfont{a}$, o teorema de Cantor
(\olref[sfr][siz][car]{thm:cantor}) e
\olref[cardinals][cardsasords]{lem:CardinalsExist} implicam que
$\cardfont{a} < \card{\Pow{\cardfont{a}}}$.
\end{proof}

\begin{thm}
O conjunto de todos os cardinais não existe.
\end{thm}

\begin{proof}
Por redução ao absurdo, suponha-se $C = \Setabs{\cardfont{a}}{\cardfont{a} \text{
é um cardinal}}$. Ora, $\bigcup C$ é um cardinal por
\olref{unioncardinalscardinal}, logo por \olref{lem:NoLargestCardinal}
existe um cardinal $\cardfont{b} > \bigcup C$. Por definição
$\cardfont{b} \in C$, logo $\cardfont{b} \subseteq \bigcup{C}$, donde
$\cardfont{b} \leq \bigcup C$, uma contradição.
\end{proof}

Compare isto com o paradoxo de Russell e com Burali-Forti.

\end{document}
